\chapter{Słaba zwartość: charakterystyki kombinatoryczne i logiczne}
\label{ch:weak-compactness-characterizations}
\chaptermark{Słaba zwartość i logika infinitarna}

\section{Trzy postacie jednej własności}

W Rozdziale~\ref{ch:mahlos-weak-compactness} wprowadziliśmy
własność drzewa oraz relacje podziałowe. Wykazaliśmy tam,
że mierzalna liczba kardynalna jest słabo zwarta
i spełnia $\kappa\longrightarrow(\kappa)^2_2$.
W tym rozdziale udowodnimy odwrotną zależność
\emph{na poziomie silnie nieosiągalnych liczb kardynalnych}
oraz powiążemy obie własności ze zwartością
infinitarnej logiki zdaniowej.

W całym rozdziale pracujemy w ZFC.
Niech $\kappa$ będzie ustaloną nieprzeliczalną
silnie nieosiągalną liczbą kardynalną.
Oznacza to, że $\kappa$ jest regularna
i $2^\lambda<\kappa$ dla każdej liczby
kardynalnej $\lambda<\kappa$.
Nie przyjmujemy, że liczba kardynalna
$\kappa$ ma własność drzewa, chyba że
występuje ona w hipotezie konkretnego twierdzenia.

\begin{twierdzenie}[Charakterystyki słabej zwartości]
\label{thm:lc39-main-equivalence}
Niech $\kappa$ będzie nieprzeliczalną
silnie nieosiągalną liczbą kardynalną.
Następujące warunki są równoważne:
\begin{enumerate}
\item $\kappa$ ma własność drzewa:
każde drzewo $\kappa$-owe ma gałąź współkońcową;
\item $\kappa\longrightarrow(\kappa)^2_2$;
\item dla każdej niezerowej liczby kardynalnej
$\mu<\kappa$ zachodzi
$\kappa\longrightarrow(\kappa)^2_\mu$;
\item każda teoria infinitarnej logiki
zdaniowej o co najwyżej $\kappa$ formułach,
z których każda zależy od mniej niż
$\kappa$ zmiennych, mająca model dla
każdej podteorii mocy mniejszej niż
$\kappa$, ma model całej teorii.
\end{enumerate}
Pojęcia z punktu czwartego zostaną dokładnie
zdefiniowane w
Definicji~\ref{def:lc39-propositional-infinitary}
oraz Definicji~\ref{def:lc39-prop-compactness}.
Dowody wszystkich implikacji podajemy poniżej.
\end{twierdzenie}

\section{Oszacowania mocy i drzewa profili kolorowania}

\begin{lemat}[Małość potęg poniżej nieosiągalności]
\label{lem:lc39-small-exponent}
Niech $\kappa$ będzie silnie nieosiągalną
liczbą kardynalną. Dla każdej liczby
kardynalnej $0<\mu<\kappa$ i każdej
liczby porządkowej $\alpha<\kappa$
zachodzi $\mu^{|\alpha|}<\kappa$.
\end{lemat}
\begin{proof}
Jeżeli $\mu=1$ albo $\alpha=0$,
lewa strona jest równa $1$.
W przeciwnym razie istnieje iniekcja
$\mu\hookrightarrow\mathcal P(\mu)$,
więc funkcję $\alpha\to\mu$ można
kodować funkcją
$\alpha\to\mathcal P(\mu)$,
a następnie podzbiorem $\alpha\times\mu$.
Otrzymujemy
\[
 \mu^{|\alpha|}
 \le 2^{|\alpha|\cdot\mu}<\kappa.
\]
Druga nierówność wynika z silnej
graniczności $\kappa$, ponieważ
$|\alpha|\cdot\mu<\kappa$:
iloczyn dwóch nieskończonych liczb
kardynalnych jest ich maksimum,
a przypadki skończone są prostsze.
\end{proof}

\begin{definicja}[Drzewo profili]
\label{def:lc39-profile-tree}
Niech $0<\mu<\kappa$ będzie liczbą
kardynalną i niech
$c:[\kappa]^2\to\mu$ będzie kolorowaniem.
Dla każdej liczby porządkowej
$\alpha<\kappa$ definiujemy
$f_\alpha:\alpha\to\mu$ przez
$f_\alpha(\xi)=c(\{\xi,\alpha\})$
dla $\xi<\alpha$.
Niech
\[
 T_c=\{f_\alpha\restriction\delta:
                   \delta\le\alpha<\kappa\},
\]
uporządkowany przez właściwe
przedłużanie funkcji.
\emph{Poziomem profili} o długości
$\delta$ jest zbiór
$T_{c,\delta}=T_c\cap\mu^\delta$.
\end{definicja}

\begin{lemat}[Drzewo profili jest drzewem $\kappa$-owym]
\label{lem:lc39-profile-kappa-tree}
Niech $\kappa$ będzie silnie nieosiągalną
liczbą kardynalną, $0<\mu<\kappa$
i $c:[\kappa]^2\to\mu$.
Drzewo $T_c$ z
Definicji~\ref{def:lc39-profile-tree}
ma wysokość $\kappa$, a
$|T_{c,\delta}|<\kappa$
dla wszystkich $\delta<\kappa$.
\end{lemat}
\begin{proof}
Zbiór $T_c$ jest zamknięty na
ograniczanie funkcji: jeśli
$t=f_\alpha\restriction\delta$,
to dla każdego $\gamma<\delta$
mamy
$t\restriction\gamma
 =f_\alpha\restriction\gamma\in T_c$.
Wysokość węzła wynosi długość
jego dziedziny. Dla każdego
$\delta<\kappa$ istnieje
$f_\alpha\restriction\delta$
z $\alpha\ge\delta$, więc
poziom $\delta$ jest niepusty.
Wysokość drzewa to $\kappa$.

Ponadto
$T_{c,\delta}\subseteq\mu^\delta$,
a z Lematu~\ref{lem:lc39-small-exponent}
wynika
$|T_{c,\delta}|
 \le\mu^{|\delta|}<\kappa$.
\end{proof}

\begin{twierdzenie}[Własność drzewa daje wszystkie małe kolorowania par]
\label{thm:lc39-tree-to-partition}
Niech $\kappa$ będzie nieprzeliczalną
silnie nieosiągalną liczbą kardynalną
mającą własność drzewa.
Dla każdej niezerowej liczby
kardynalnej $\mu<\kappa$ zachodzi
\[
 \kappa\longrightarrow(\kappa)^2_\mu.
\]
\end{twierdzenie}
\begin{proof}
Ustalmy $0<\mu<\kappa$ i
$c:[\kappa]^2\to\mu$.
Z Lematu~\ref{lem:lc39-profile-kappa-tree}
drzewo profili $T_c$ jest
drzewem $\kappa$-owym.
Na mocy własności drzewa
ma gałąź współkońcową.
Jej suma jest funkcją
$b:\kappa\to\mu$;
dla każdego $\delta<\kappa$
jej ograniczenie
$b\restriction\delta$
należy do $T_{c,\delta}$.

Dla każdego $\delta<\kappa$
istnieje $\alpha\ge\delta$
taka, że
$f_\alpha\restriction\delta
 =b\restriction\delta$.
Wybieramy \emph{najmniejszą}
taką liczbę porządkową
i oznaczamy ją $a(\delta)$.
Zatem
\[
 a(\delta)\ge\delta,\qquad
 f_{a(\delta)}\restriction\delta
       =b\restriction\delta.
\]
Definiujemy $h:\kappa\to\mu$,
$h(\delta)=b(a(\delta))$.

Ponieważ $\kappa$ jest regularna,
pewien kolor $i<\mu$ ma
\emph{nieograniczony} przeciwobraz
$D=\{\delta<\kappa:h(\delta)=i\}$.
W istocie, gdyby każdy przeciwobraz
był ograniczony, supremum
mniej niż $\kappa$ ich ograniczeń
pozostałoby mniejszy od $\kappa$,
co przeczyłoby temu, że
ich suma zbiorów jest równa $\kappa$.

Rekurencyjnie, dla $\xi<\kappa$,
wybieramy $\delta_\xi\in D$
tak, by
\[
 \delta_\xi>
    \sup_{\eta<\xi}(a(\delta_\eta)+1).
\]
Istnieje odpowiedni punkt,
gdyż $\xi<\kappa$,
regularność zapewnia ograniczoność
wyrażenia po prawej, a $D$
jest nieograniczony.
Możemy wybierać najmniejszy
punkt spełniający wymaganie.
Niech $\alpha_\xi=a(\delta_\xi)$.
Z konstrukcji mamy
$\alpha_\eta<\delta_\xi\le\alpha_\xi$
dla każdej pary $\eta<\xi$,
więc $(\alpha_\xi)_{\xi<\kappa}$
jest ściśle rosnący.

Dla $\eta<\xi$:
\[
 \begin{aligned}
 c(\{\alpha_\eta,\alpha_\xi\})
 &=f_{\alpha_\xi}(\alpha_\eta)\\
 &=b(\alpha_\eta)
   =h(\delta_\eta)=i.
 \end{aligned}
\]
Pierwsza równość jest definicją
$f_{\alpha_\xi}$, druga wynika z
$\alpha_\eta<\delta_\xi$ i
$f_{\alpha_\xi}\restriction\delta_\xi
 =b\restriction\delta_\xi$,
a ostatnia z $\delta_\eta\in D$.
Zatem $H=\{\alpha_\xi:\xi<\kappa\}$
ma moc $\kappa$ i wszystkie
jego pary otrzymują kolor $i$.
\end{proof}

\section{Własność podziałowa wymusza gałąź}

\begin{definicja}[Porządek leksykograficzny drzewa]
\label{def:lc39-tree-lex}
Niech $(T,<_{T})$ będzie drzewem
o wysokości $\kappa$.
Dla każdego $\delta<\kappa$
wybieramy porządek liniowy
$\prec_\delta$ na poziomie $T_\delta$.

Każdemu węzłowi $t$ przyporządkowujemy
\emph{ścieżkę poprzedników}
$p_t$, której dziedziną jest
$\operatorname{ht}_{T}(t)+1$ i
$p_t(\delta)$ jest jedynym
węzłem na poziomie $\delta$
nie większym od $t$.
Na $T$ definiujemy $\prec$
jako porządek leksykograficzny
tych ścieżek: ścieżka będąca
właściwym początkiem drugiej
jest wcześniejsza; w przeciwnym
razie porównujemy węzły
w pierwszym poziomie $\delta$,
na którym ścieżki się różnią,
korzystając z $\prec_\delta$.
\end{definicja}

\begin{lemat}[Rozdzielenie stożków w porządku leksykograficznym]
\label{lem:lc39-lex-cones}
Niech $\prec$ będzie porządkiem
z Definicji~\ref{def:lc39-tree-lex}.
Jest to porządek liniowy
przedłużający porządek $<_{T}$.
Jeśli węzły $s,t\in T$ są
nieporównywalne w $T$ i
$s\prec t$, to dla wszystkich
$s'\ge_T s$ i $t'\ge_T t$
mamy $s'\prec t'$.
\end{lemat}
\begin{proof}
Dwie różne ścieżki poprzedników
są porównywalne leksykograficznie:
jeżeli żadna nie jest początkiem
drugiej, istnieje najmniejszy
poziom, na którym się różnią.
Porządki $\prec_\delta$ są liniowe,
więc opisane porównanie jest
jednoznaczne. Przechodniość
wynika z rozważenia pierwszego
poziomu różnicy trzech ścieżek;
daje to standardowy dowód
przechodniości porządku
leksykograficznego.
Jeżeli $s<_{T}t$, to
$p_s$ jest właściwym początkiem
$p_t$, więc $s\prec t$.

Załóżmy, że $s,t$ są
nieporównywalne.
Ich ścieżki różnią się
na pewnym najmniejszym
poziomie $\delta$.
Wszystkie przedłużenia $s'$
i $t'$ mają na poziomach
$\le\delta$ te same węzły,
co odpowiednio $s$ i $t$.
Zatem pierwsza różnica
między ścieżkami $s',t'$
występuje na tym samym
poziomie $\delta$ i ma
ten sam znak w
porządku $\prec_\delta$.
Stąd $s'\prec t'$.
\end{proof}

\begin{twierdzenie}[Relacja podziałowa implikuje własność drzewa]
\label{thm:lc39-partition-to-tree}
Niech $\kappa$ będzie nieprzeliczalną
regularną liczbą kardynalną
spełniającą
$\kappa\longrightarrow(\kappa)^2_2$.
Wtedy każde drzewo wysokości
$\kappa$, którego poziomy mają
moc mniejszą od $\kappa$,
ma gałąź współkońcową.
\end{twierdzenie}
\begin{proof}
Niech $T$ będzie dowolnym
takim drzewem. Ponieważ
ma wysokość $\kappa$,
każdy poziom $\delta<\kappa$
jest niepusty. Z jednej strony
$|T|\ge\kappa$, z drugiej
$|T|\le\kappa\cdot\kappa=\kappa$,
więc $|T|=\kappa$.
W ZFC wybieramy wyliczenie
$T=\{t_\alpha:\alpha<\kappa\}$
bez powtórzeń i porządek
leksykograficzny $\prec$
z Definicji~\ref{def:lc39-tree-lex}.

Definiujemy kolorowanie
$c:[\kappa]^2\to\{0,1\}$:
dla $\alpha<\beta$ kładziemy
\[
 c(\{\alpha,\beta\})=
 \begin{cases}
 0,&t_\alpha\prec t_\beta,\\
 1,&t_\beta\prec t_\alpha.
 \end{cases}
\]
Z założenia istnieje
$H\subseteq\kappa$ o mocy
$\kappa$, na którym $c$
jest stale równe $i<2$.
Dla $s\in T$ określamy
\[
 U_s=\{\alpha\in H:s\le_T t_\alpha\},
 \qquad
 A=\{s\in T:|U_s|=\kappa\}.
\]

Pokażemy, że dla każdego
$\delta<\kappa$ zbiór
$A\cap T_\delta$ jest niepusty.
Ponieważ $\kappa$ jest regularna,
zbiór wszystkich węzłów z
poziomów mniejszych od $\delta$
ma moc mniejszą od $\kappa$:
jest sumą zbiorów mniej niż
$\kappa$ poziomów, każdy
o mocy mniejszej od $\kappa$.
Zatem dla $\kappa$ wielu
indeksów $\alpha\in H$
węzeł $t_\alpha$ ma
wysokość co najmniej $\delta$.
Każdy taki węzeł ma jednego
poprzednika na poziomie $\delta$.
Poziom ten ma moc mniejszą
od $\kappa$; regularność
gwarantuje, że dla co najmniej
$\kappa$ wielu rozważanych
indeksów poprzednik jest
tym samym węzłem $s$.
Wtedy $s\in A\cap T_\delta$.

Pozostaje pokazać, że
węzły $A$ są parami
porównywalne w $T$.
Przypuśćmy przeciwnie,
że $s,t\in A$ są
nieporównywalne i,
po zamianie oznaczeń,
$s\prec t$.
Z Lematu~\ref{lem:lc39-lex-cones}
dla każdych
$\alpha\in U_s$ i
$\beta\in U_t$ mamy
$t_\alpha\prec t_\beta$.
Zbiory $U_s$ i $U_t$
są rozłączne: jeden węzeł
nie może przedłużać dwóch
nieporównywalnych węzłów.

Jeśli $i=0$, to jednorodność
$H$ oznacza, że
$\alpha<\beta$ zawsze, gdy
$\alpha\in U_s$ i
$\beta\in U_t$.
Jednak $|U_s|=\kappa$
oznacza, że $U_s$ jest
nieograniczony w $\kappa$;
wybierając dowolne
$\beta\in U_t$ i większe
$\alpha\in U_s$, otrzymujemy
sprzeczność.
Jeśli $i=1$, jednorodność
oznacza $\beta<\alpha$
dla wszystkich takich par.
Wybierając dowolne
$\alpha\in U_s$ i większe
$\beta\in U_t$ ponownie
otrzymujemy sprzeczność.
Zatem $A$ jest łańcuchem.
Przecina wszystkie poziomy
$T_\delta$, więc stanowi
gałąź współkońcową.
\end{proof}

\begin{wniosek}[Równoważność dwóch charakterystyk]
\label{cor:lc39-tree-partition}
Dla każdej nieprzeliczalnej
silnie nieosiągalnej liczby
kardynalnej $\kappa$
następujące warunki są
równoważne:
\[
 \kappa\text{ ma własność drzewa}
 \quad\Longleftrightarrow\quad
 \kappa\longrightarrow(\kappa)^2_2.
\]
\end{wniosek}
\begin{proof}
Implikacja z lewej na prawą
jest przypadkiem $\mu=2$
Twierdzenia~\ref{thm:lc39-tree-to-partition}.
Implikacja odwrotna wynika
z Twierdzenia~\ref{thm:lc39-partition-to-tree}.
\end{proof}

\begin{uwaga}[Co zawiera hipoteza nieosiągalności?]
\label{rem:lc39-scope}
W tym rozdziale dowodzimy
równoważności \emph{przy założeniu}
silnej nieosiągalności.
W pełnej charakterystyce
nieprzeliczalnych słabo zwartych
liczb kardynalnych sam warunek
$\kappa\to(\kappa)^2_2$
implikuje również silną
nieosiągalność. Dowód tego
dodatkowego faktu wymaga
osobnego argumentu o porządku
leksykograficznym na $2^\lambda$;
nie używamy go tutaj bez dowodu.
\end{uwaga}

\section{Zwartość infinitarnej logiki zdaniowej}

\begin{definicja}[Infinitarna logika zdaniowa o małych nośnikach]
\label{def:lc39-propositional-infinitary}
Niech $\kappa$ będzie nieprzeliczalną
regularną liczbą kardynalną.
Dla każdego $\xi<\kappa$
niech $p_\xi$ będzie odrębną
zmienną zdaniową.
Rozważamy formuły utworzone
z tych zmiennych przez
negację oraz skończone
lub nieskończone koniunkcje
i alternatywy indeksowane
zbiorami o mocy mniejszej
niż $\kappa$.
Zbiór zmiennych występujących
w danej formule $\varphi$
nazywamy jej \emph{nośnikiem}
i oznaczamy $\operatorname{supp}(\varphi)$.
Dopuszczamy w naszej teorii
wyłącznie formuły o nośniku
mocy mniejszej niż $\kappa$.

\emph{Wartościowanie} to funkcja
$v:\kappa\to\{0,1\}$.
Relację $v\models\varphi$
określamy rekurencyjnie:
$v\models p_\xi$ wtedy i tylko
wtedy, gdy $v(\xi)=1$;
negacja zmienia wartość
logiczną, koniunkcja
wymaga prawdziwości
wszystkich składników,
a alternatywa co najmniej
jednego ze składników.
Z definicji
$v\models\varphi$ zależy
wyłącznie od ograniczenia
$v\restriction\operatorname{supp}(\varphi)$.
\end{definicja}

\begin{definicja}[Lokalna spełnialność i zwartość]
\label{def:lc39-prop-compactness}
Niech $\kappa$ będzie nieprzeliczalną
regularną liczbą kardynalną,
a $\Sigma$ zbiorem co najwyżej
$\kappa$ formuł z
Definicji~\ref{def:lc39-propositional-infinitary}.
Mówimy, że $\Sigma$ jest
\emph{$\kappa$-lokalnie spełnialna},
gdy dla każdego
$\Sigma_0\subseteq\Sigma$
o mocy $|\Sigma_0|<\kappa$
istnieje wartościowanie
$v:\kappa\to\{0,1\}$
spełniające wszystkie
formuły z $\Sigma_0$.

Mówimy, że zachodzi
\emph{zasada $\kappa$-zwartości
infinitarnej logiki zdaniowej},
gdy każda taka
$\kappa$-lokalnie spełnialna
$\Sigma$ ma wartościowanie
spełniające wszystkie
jej formuły.
\end{definicja}

\begin{lemat}[Drzewo częściowych wartościowań]
\label{lem:lc39-valuations-tree}
Niech $\kappa$ będzie nieprzeliczalną
silnie nieosiągalną liczbą
kardynalną, a
$\Sigma=\{\varphi_\xi:\xi<\kappa\}$
niech będzie
$\kappa$-lokalnie spełnialną
teorią z Definicji~\ref{def:lc39-prop-compactness}.
Dla $\alpha<\kappa$ oznaczmy
\[
 \begin{aligned}
 T_\alpha=\{v\in2^\alpha:\ 
 &v\text{ spełnia }\varphi_\xi\
   \text{ zawsze, gdy }\xi<\alpha\\
 &\text{oraz }
   \operatorname{supp}(\varphi_\xi)
                           \subseteq\alpha\}.
 \end{aligned}
\]
Przez spełnianie formuły
przez częściowe $v$
rozumiemy jej obliczenie
na wszystkich zmiennych
w jej nośniku.
Wtedy
$T=\bigcup_{\alpha<\kappa}T_\alpha$
uporządkowany przez
przedłużanie funkcji
jest drzewem $\kappa$-owym.
\end{lemat}
\begin{proof}
Dla ustalonego $\alpha$
formuł o indeksach $\xi<\alpha$,
których nośniki zawierają się
w $\alpha$, jest mniej niż $\kappa$.
Z lokalnej spełnialności
istnieje pełne wartościowanie
spełniające te formuły.
Jego ograniczenie do $\alpha$
należy do $T_\alpha$,
więc każdy poziom jest niepusty.

Jeśli $v\in T_\alpha$
oraz $\beta<\alpha$,
to $v\restriction\beta$
spełnia wszystkie
$\varphi_\xi$ z $\xi<\beta$
o nośniku zawartym w $\beta$:
były one również warunkami
dla $v$ na poziomie $\alpha$.
Zatem $v\restriction\beta
\in T_\beta$.
Zbiór $T$ jest więc drzewem
wysokości $\kappa$ z
poziomami $T_\alpha$.
Ponadto
$|T_\alpha|
 \le 2^{|\alpha|}<\kappa$
na mocy silnej graniczności.
\end{proof}

\begin{twierdzenie}[Własność drzewa daje zwartość zdaniową]
\label{thm:lc39-tree-to-compactness}
Niech $\kappa$ będzie nieprzeliczalną
silnie nieosiągalną liczbą
kardynalną mającą własność drzewa.
Wtedy zachodzi zasada
$\kappa$-zwartości infinitarnej
logiki zdaniowej z
Definicji~\ref{def:lc39-prop-compactness}.
\end{twierdzenie}
\begin{proof}
Ustalmy $\kappa$-lokalnie
spełnialną teorię $\Sigma$
o mocy co najwyżej $\kappa$.
Jeżeli $|\Sigma|<\kappa$,
jej spełnialność jest
bezpośrednio częścią hipotezy.
Załóżmy więc $|\Sigma|=\kappa$
i wybierzmy wyliczenie
$\Sigma=\{\varphi_\xi:\xi<\kappa\}$.

Na mocy
Lematu~\ref{lem:lc39-valuations-tree}
drzewo częściowych wartościowań
jest drzewem $\kappa$-owym.
Z własności drzewa ma gałąź
współkońcową. Jej suma
jest pełnym wartościowaniem
$v:\kappa\to\{0,1\}$,
bo gałąź przechodzi przez
każdy poziom $\alpha<\kappa$.

Ustalmy dowolne $\xi<\kappa$.
Z założenia
$|\operatorname{supp}(\varphi_\xi)|
 <\kappa$ i z regularności
$\kappa$ wynika, że
$\operatorname{supp}(\varphi_\xi)$
jest ograniczony w $\kappa$.
Możemy więc wybrać
$\alpha<\kappa$ większe od
$\xi$ i wszystkich
indeksów z tego nośnika.
Węzeł $v\restriction\alpha$
należy do $T_\alpha$,
a zatem spełnia
$\varphi_\xi$.
Ponieważ prawdziwość
zależy tylko od nośnika,
pełne $v$ również spełnia
$\varphi_\xi$.
Formuła była dowolna,
więc $v\models\Sigma$.
\end{proof}

\begin{twierdzenie}[Zwartość zdaniowa daje własność drzewa]
\label{thm:lc39-compactness-to-tree}
Niech $\kappa$ będzie
nieprzeliczalną regularną
liczbą kardynalną.
Załóżmy zasadę
$\kappa$-zwartości
infinitarnej logiki zdaniowej
z Definicji~\ref{def:lc39-prop-compactness}.
Wtedy każde drzewo
$\kappa$-owe ma gałąź
współkońcową.
\end{twierdzenie}
\begin{proof}
Niech $T$ będzie drzewem
$\kappa$-owym.
Ma ono dokładnie
$\kappa$ węzłów,
ponieważ ma wysokość $\kappa$
i mniej niż $\kappa$ węzłów
na każdym z $\kappa$ poziomów.
Wybieramy dla każdego
$t\in T$ odrębną zmienną
zdaniową $p_t$.
Po wyliczeniu węzłów
możemy utożsamić te zmienne
z rodziną $(p_\xi)_{\xi<\kappa}$.

Tworzymy teorię $\Sigma_T$
z dwóch rodzajów formuł.
Dla każdego poziomu
$\alpha<\kappa$ dodajemy
\[
 P_\alpha=
 \bigvee_{t\in T_\alpha}p_t.
\]
Jest to dopuszczalna formuła
infinitarna, ponieważ
$|T_\alpha|<\kappa$.
Dla każdej pary
nieporównywalnych węzłów
$s,t\in T$ dodajemy
\[
 Q_{s,t}=\neg(p_s\land p_t).
\]
Cała teoria ma moc
co najwyżej $\kappa$,
ponieważ liczba poziomów
wynosi $\kappa$,
a liczba par węzłów
jest co najwyżej
$\kappa^2=\kappa$.
Każda formuła zależy od
mniej niż $\kappa$ zmiennych.

Wykażemy lokalną spełnialność.
Niech $\Sigma_0\subseteq\Sigma_T$
ma moc mniejszą od $\kappa$.
Zbiór poziomów $\alpha$,
dla których
$P_\alpha\in\Sigma_0$,
ma moc mniejszą od $\kappa$.
Ponieważ $\kappa$ jest regularna,
istnieje $\delta<\kappa$
większa od wszystkich
tych poziomów.
Wybieramy dowolny
$t\in T_\delta$.
Definiujemy wartościowanie
\[
 v_t(p_s)=1
   \quad\Longleftrightarrow\quad s\le_T t.
\]
Wszystkie węzły na ścieżce
poprzedników $t$ są
porównywalne, więc
$v_t$ spełnia wszystkie
formuły typu $Q_{s,u}$.
Dla każdego poziomu
$\alpha$ uwzględnionego
w $\Sigma_0$ węzeł $t$
ma poprzednika na poziomie
$\alpha$, a więc
$v_t$ spełnia także
$P_\alpha$.
W konsekwencji
$v_t\models\Sigma_0$.

Z założonej zwartości
otrzymujemy wartościowanie
$v\models\Sigma_T$.
Dla każdego $\alpha<\kappa$
formuła $P_\alpha$ wymusza,
że istnieje węzeł
$t_\alpha\in T_\alpha$
z $v(p_{t_\alpha})=1$.
Wszystkie wybrane węzły
muszą być parami
porównywalne, ponieważ
dla nieporównywalnej pary
odpowiednia formuła
$Q_{s,t}$ zakazuje
jednoczesnej wartości $1$.
Zbiór węzłów o wartości $1$
jest zatem łańcuchem
przecinającym każdy poziom,
a więc gałęzią
współkońcową.
\end{proof}

\begin{wniosek}[Kombinatoryczna i zdaniowa zwartość]
\label{cor:lc39-compactness-equivalence}
Niech $\kappa$ będzie nieprzeliczalną
silnie nieosiągalną liczbą
kardynalną. Wtedy
własność drzewa, relacja
$\kappa\longrightarrow(\kappa)^2_2$
oraz zasada $\kappa$-zwartości
infinitarnej logiki zdaniowej
są równoważne.
\end{wniosek}
\begin{proof}
Twierdzenia~\ref{thm:lc39-tree-to-compactness}
i~\ref{thm:lc39-compactness-to-tree}
dają równoważność własności
drzewa i zwartości zdaniowej.
Wniosek~\ref{cor:lc39-tree-partition}
daje równoważność
własności drzewa i relacji
podziałowej.
\end{proof}

\begin{przyklad}[Teoria opisująca gałąź w skończonym drzewie]
\label{ex:lc39-finite-branch}
Rozważmy drzewo z trzema poziomami
\[
 \{\varnothing\},\quad
 \{(0),(1)\},\quad
 \{(0,0),(0,1),(1,0),(1,1)\},
\]
uporządkowane przez
przedłużanie skończonych ciągów.
Wprowadzamy zmienną $p_s$
dla każdego węzła $s$.
Formuła pierwszego poziomu
to $p_\varnothing$.
Drugi poziom daje
$p_{(0)}\lor p_{(1)}$,
a trzeci alternatywę
czterech odpowiednich zmiennych.
Dodatkowo wymagamy
$\neg(p_{(0)}\land p_{(1)})$
oraz analogicznych
ograniczeń dla
wszystkich nieporównywalnych
par węzłów.

Wartościowanie przyjmujące
wartość $1$ dokładnie
na węzłach
$\varnothing,(0),(0,1)$
spełnia wszystkie formuły:
na każdym poziomie
wybiera co najmniej jeden
węzeł, a wybrane
węzły są porównywalne.
Tak działa kodowanie
gałęzi w dowodzie
Twierdzenia~\ref{thm:lc39-compactness-to-tree}.
\end{przyklad}

\section{Dodatkowa konieczna własność regularności}

\begin{twierdzenie}[Relacja podziałowa wyklucza singularność]
\label{thm:lc39-partition-regular}
Niech $\kappa$ będzie
nieprzeliczalną liczbą kardynalną.
Jeśli
$\kappa\longrightarrow(\kappa)^2_2$,
to $\kappa$ jest regularna.
\end{twierdzenie}
\begin{proof}
Przypuśćmy, że
$\operatorname{cf}(\kappa)=\mu<\kappa$.
Wybieramy rosnący
ciąg $(\gamma_\xi)_{\xi<\mu}$
współkońcowy w $\kappa$
i dzielimy $\kappa$
na $\mu$ rozłącznych
przedziałów $I_\xi$,
z których każdy ma
moc mniejszą od $\kappa$.
Można je otrzymać,
biorąc przedziały między
kolejnymi punktami ciągu;
występujące ewentualnie
puste przedziały nie
zmieniają argumentu.

Kolorujemy parę
$\{\alpha,\beta\}$ kolorem
$0$, gdy oba elementy
należą do jednego $I_\xi$,
oraz kolorem $1$
w przeciwnym wypadku.
Jeżeli $H\subseteq\kappa$
jest jednorodny
w kolorze $0$, to
wszystkie jego elementy
należą do jednego
przedziału, zatem
$|H|<\kappa$.
Jeśli $H$ jest
jednorodny w kolorze $1$,
to zawiera co najwyżej
jeden punkt każdego
przedziału, a więc
$|H|\le\mu<\kappa$.
Nie istnieje zatem
jednorodny zbiór
mocy $\kappa$;
sprzeczność z hipotezą.
\end{proof}

\section{Podsumowanie i zakres dalszego rozwoju}

\begin{proof}[Dowód Twierdzenia~\ref{thm:lc39-main-equivalence}]
Implikacja (1)$\Rightarrow$(3)
została udowodniona w
Twierdzeniu~\ref{thm:lc39-tree-to-partition}.
Implikacja (3)$\Rightarrow$(2)
wynika przez podstawienie
$\mu=2$.
Implikacja (2)$\Rightarrow$(1)
jest Twierdzeniem~\ref{thm:lc39-partition-to-tree}.
Równoważność punktów (1) i (4)
zawierają
Twierdzenia~\ref{thm:lc39-tree-to-compactness}
oraz~\ref{thm:lc39-compactness-to-tree}.
Wszystkie warunki są więc
równoważne.
\end{proof}

\begin{uwaga}[Zwartość zdaniowa i logika pierwszego rzędu]
\label{rem:lc39-first-order}
Zasada zwartości użyta w tym
rozdziale dotyczy
\emph{infinitarnej logiki zdaniowej}
z formułami zależnymi od
mniej niż $\kappa$ zmiennych.
Nie utożsamiamy jej bez dowodu
z pełnym twierdzeniem
o zwartości logiki pierwszego
rzędu $L_{\kappa,\kappa}$,
dopuszczającej również bloki
kwantyfikatorów długości
mniejszej niż $\kappa$.
Pełne przejście między tymi
językami wymaga osobnego
dowodu konstrukcji modelu
i odpowiedniej wersji
rachunku Henkina.
\end{uwaga}

\begin{uwaga}[Zależność od silnej nieosiągalności]
\label{rem:lc39-inaccessible-usage}
W dowodzie z własności
drzewa do relacji podziałowej
kluczowe było oszacowanie
$\mu^{|\delta|}<\kappa$,
które zapewnia małe poziomy
drzewa profili.
Z kolei przy zwartości
zdaniowej małe poziomy
drzewa częściowych
wartościowań daje
$2^{|\delta|}<\kappa$.
Nie można pominąć
tych oszacowań i twierdzić,
że samo istnienie
drzewa o wysokości
$\kappa$ wystarcza.
\end{uwaga}

\begin{uwaga}[Status formalnej weryfikacji]
\label{rem:lc39-formal}
Przedstawione dowody są
dowodami matematycznymi
w tekście LaTeX.
Automatyczna kontrola
odnośników, środowisk
i skończonych testów
kombinatorycznych nie
zastępuje formalizacji
w Lean. Status formalizacji
wymaga odrębnego,
rzeczywistego przebiegu
kompilatora Lean.
\end{uwaga}
